\documentstyle[% 


righttag% 


,11pt% 


,amssymb% 


,russian%               remove in Eng. 


,izvru%                 Set izven% in Eng. 


]{amsart} 


 


%  {=}


%% !!! \varphi'_\wedge, \varphi'_\sqcap = левая и правая производные.


%% \wedge и \sqcap выглядят как Л и П. Как менять???????????


%% На \mathrm L и \mathrm R ???????


\newcommand{\const}{\operatorname{const}} 


\newcommand{\sign}{\operatorname{sign}} 


\newcommand{\tts}[1]{} 


 


\theoremstyle{plain} 


\theorembodyfont{\it} 


\newtheorem{assrtnonum}{\hskip\parindent Утверждение} 


\renewcommand{\theassrtnonum}{} 


 


 


%\hoffset=-15mm%                  For Eng. 


\setcounter{page}{49} 


\god{2005} 


\nomer{8~(519)} %                        Remove in Eng. 


%\nomer{Vol.\,49, No.\,8}%               Set in Eng. 


%\str{\pageref{begin}--\pageref{end}}%  Set in Eng. 


\udk{517.518} 


 


\author{\it Б.В.\,Симонов} 


% NB:   ^^ remove in Eng. 


\title{Несимметричные приближения функций 


многих переменных в функциональных пространствах} 


 


\date{} 


%\pagestyle{myheadings}%         Set in Eng. 


\pagestyle{plain} %              Remove in Eng. 


\begin{document} 


%\thispagestyle{plain}%          Set in Eng. 


\maketitle 


\label{begin} 


 


 


 


Задача об элементе наилучшего приближения в различных функциональных 


пространствах является одной из основных задач теории приближений. В 


данной работе рассматриваются несимметричные приближения для функций 


нескольких переменных в пространствах, обобщающих пространства $L_r$. 


Введем ряд определений. 


 


Пусть $\phi$ --- совокупность неотрицательных, неубывающих, непрерывных и 


выпуклых вниз на полупрямой $[0,+\infty)$ функций $\varphi$, 


удовлетворяющих $\Delta_2$-условию (т.\,е. существует такая положительная 


постоянная $c$, что $\varphi(2u)\le c\varphi(u)$ для всех 


$u\in(0,+\infty)$). Пусть $G$ --- измеримое множество в $R^n$, где $R^n$ 


--- $n$-мерное евклидово пространство точек $x=(x_1,x_2,\dotsc,x_n)$ с 


действительными координатами. Через $L_{\varphi,G}$, $\varphi\in\phi$, 


обозначим множество всех измеримых, действительных функций $f(x)$, 


заданных на $G$ и таких, что 


$\|f\|_{\varphi,G}=\int\limits_G\varphi(|f(x)|)dx<+\infty$. 


 


Пусть $L_0(G)$ --- линейное пространство всех измеримых, конечных 


почти всюду (п.\,в.) на $G$ функций. Считаем $f_1=f_2$, если 


$$ 


\mu\{x\in G:f_1(x)\ne f_2(x)\}=0. 


$$ 


 


Упорядоченную пару $p=(p_1,p_2)$, $p_1,p_2\in L_0(G)$, назовем весом на 


$G$. Введем обозначение $(f;p)(x)=f^+(x)p_1(x)-f^-(x)p_2(x)$, где 


$f^+=\max\{f,0\}$, $f^-=\max\{-f,0\}$. Заметим, что если $p_1(x)\equiv 


p_2(x)\equiv1$, то $(f;p)(x)=f(x)$. Рассмотрим функционал 


$\|(f;p)\|_{\varphi,G}$. Если $\varphi(u)=u^r$ ($r\ge1$), $n=1$, 


$G=[a,b]$, $p_1(x)\equiv\alpha>0$, $p_2(x)\equiv\beta>0$, то 


$\|(f;p)\|^{1/r}_{\varphi,G}$ совпадает с функционалом 


$\|f\|_{r,(\alpha,\beta)}$, введенным в [1]. 


 


Пусть $H$ --- конечномерное подпространство пространства $L_{\varphi,G}$, 


$f,g\in L_{\varphi,G}$. Положим $\rho(f,g)=\|(f-g;p)\|_{\varphi,G}$, 


\begin{equation} 


\rho(f,H)=\inf_{g\in H}\rho(f,g). 


\end{equation} 


 


Величину $\rho(f,H)$ назовем наилучшим несимметричным приближением 


функции $f$ элементами из $H$ в пространстве $L_{\varphi,G}$, а элемент 


$g_0\in H$, на котором достигается \ $\inf$ \ в (1) --- элементом 


наилучшего несимметричного приближения в $L_{\varphi,G}$ подпространством 


$H$. Если $\varphi=u^r$ ($r\ge1$), $G=[a,b]^n$, $p_1(x)\equiv 


p_2(x)\equiv1$, то $\rho^{1/r}(f,H)$ является наилучшим приближением 


функции $f$ элементами из $H$ в метрике $L_r$. 


 


Задача о существовании и нахождении элемента наилучшего приближения в 


различных функциональных пространствах рассматривалась во многих работах 


[2]--[4]. 


 


Этим же вопросам в пространствах $L_{\varphi,G}$ посвящена данная работа. 


Основной результат статьи содержится в теоремах 1, 2. 


 


 


\begin{thm} 


Пусть $\varphi\in\phi$, $p=(p_1,p_2)$ --- такой вес, что 


\begin{equation} 


\sum^2_{i=1}\mu\{x\in G:p_i(x)=0\}=0. 


\end{equation} 


Тогда для любого конечномерного подпространства $H\in L_{\varphi,G}$ и 


любой функции $f\in L_{\varphi,G}$ существует элемент наилучшего 


несимметричного приближения в $L_{\varphi,G}$ подпространством $H$. Если 


к тому же $\varphi(\frac{u_1+u_2}{2})< 


\frac{\varphi(u_1)+\varphi(u_2)}{2}$, $0\le u_1<u_2<\infty$ $($условие 


строгой выпуклости функции~$\varphi)$, то элемент наилучшего 


несимметричного приближения будет единственным. 


\end{thm} 


 


 


\begin{thm} 


Пусть $\varphi\in \phi$, $H$ --- фиксированное конечномерное 


подпространство из $L_{\varphi,G}$, $p=(p_1,p_2)$ --- такой вес из 


$L_0(G)$, что $gp_1\in L_{\varphi,G}$, $gp_2\in L_{\varphi,G}$ 


для любого элемента $g\in H$. 


 


Тогда для того чтобы функция $g_0\in H$ доставляла функции $f\in 


L_{\varphi,G}$ наилучшее несимметричное приближение, необходимо и 


достаточно, чтобы при любой функции $g\in H$ выполнялось условие для 


$F(x)=f(x)-g_0(x)$ 


\begin{multline} 


\int_{G\setminus E}\big\{\varphi'_\wedge(|(F(x);p(x))|)\sign\{F(x)g(x) 


\}^-+\varphi'_\sqcap(|(F(x);p(x))|)\sign 


\{F(x)g(x)\}^+\big\}\times \\ 


% 


\times g(x)(|p_1(x)|\sign(F(x))^+-|p_2(x)|\sign(F(x))^-) 


dx+\varphi'_\sqcap(0)\int_E|(g(x);p(x))|dx\ge0, 


\end{multline} 


где $E=\{x\in G:F(x)=0\}$,


$\varphi'_\wedge$ --- левая производная, $\varphi'_\sqcap$ --- правая производная. 


\end{thm} 


 


 


\begin{note} 


Пусть функция $\varphi$ дифференцируема на $(0,+\infty)$. 


Тогда нетрудно показать, что в теореме 2 условие (3) можно заменить 


на эквивалентное ему условие 


\begin{multline*} 


\bigg|\int_{G\setminus E}\varphi'(|(F(x);p(x))|)g(x) 


(|p_1(x)|\sign(F(x))^+ -\\ 


% 


-|p_2(x)|\sign(F(x))^-)dx\bigg|\le\varphi'_\sqcap(0) 


\int_E|(g(x);p(x))|dx. 


\end{multline*} 


\end{note} 


 


 


\begin{note} 


Пусть $\varphi(u)=u^r$ ($r\ge1$), $G=(a_1,b_1)\times \dotsb \times 


(a_n,b_n)$. Тогда из теоремы 2 и замечания 1 нетрудно получить 


\end{note} 


 


 


\begin{assrtnonum} 


Пусть $H$ --- конечномерное подпространство в $L_r(G)$ $(1\le 


r<+\infty)$. Для того чтобы функция $g_0\in H$ доставляла функции $f\in 


L_r(G)$ наилучшее приближение в подпространстве $H$, необходимо и 


достаточно, чтобы для любой функции $g\in H$ выполнялось условие 


\begin{equation} 


\bigg|\int_{G\setminus E}g(x)|F(x)|^{r-1}\sign F(x)dx\bigg|\le 


(1-\sign(r-1))\int_E|g(x)|dx. 


\end{equation} 


\end{assrtnonum} 


 


 


\begin{note} 


Пусть $\varphi(u)=u^r$ ($r\ge1$), $n=2$, $G=[0,2\pi]^2$, $H=T_{M_1M_2}$ 


--- множество тригонометрических полиномов двух переменных порядка не 


выше $M_1$, $M_2$ по переменным $x_1$, $x_2$ соответственно. 


 


Для того чтобы тригонометрический полином $t^0_{k_1k_2}(x_1,x_2)$ из 


$T_{M_1M_2}$ доставлял функции\linebreak


$f\in L_r(0,2\pi)^2$ наилучшее приближение 


в подпространстве $T_{M_1M_2}$, необходимо и достаточно, чтобы для любого 


тригонометрического полинома $t_{m_1m_2}(x_1,x_2)$ из этого 


подпространства выполнялось условие при 


$F(x_1,x_2)=f(x_1,x_2)-t^0_{k_1k_2}(x_1,x_2)$ 


\begin{multline*} 


\bigg|\;\iint\limits_{[0,2\pi]^2\setminus E}t_{m_1m_2}(x_1,x_2) 


|F(x_1,x_2)|^{r-1}\sign F(x_1,x_2)dx_1\,dx_2\bigg|\le \\ 


% 


\le(1-\sign(r-1))\iint\limits_E|t_{m_1m_2}(x_1,x_2)|dx_1\,dx_2, 


\end{multline*} 


где $E=\{(x_1,x_2)\in[0,2\pi]^2:F(x_1,x_2)=0\}$. 


\end{note} 


 


 


\begin{note} 


Приведем пример, показывающий, что равенство в (4) достижимо. 


\end{note} 


 


 


\begin{exam} 


Пусть $r=1$, $f(x_1,x_2)=\sign x_1\sign x_2$, $G=[-1,1]^2$, $H=\{C\}$ --- 


множество постоянных. Тогда элементы наилучшего приближения образуют 


множество таких постоянных $C$, что $|C|\le1$. 


 


Если $g_0=C\in(-1,1)$, то условие (4) имеет вид $0=0$; если $g_0=-1$ 


или $g_0=1$, то получим $2|C|=2|C|$ для любого $C\in H$. 


\end{exam} 


 


 


\begin{note} 


Если в теореме 1 для веса $p=(p_1,p_2)$ не выполняется условие (2), то 


существуют пространства $L_{\varphi,G}$ и такие их подпространства $H$, в 


которых нет элемента наилучшего несимметричного приближения. Это 


иллюстрирует 


\end{note} 


 


 


\begin{exam} 


Пусть $\varphi(u){=}u$, $G{=}[0,1]^2$, $H{=}\{C\}$, $p_1(x_1,x_2){=}1$, 


$p_2(x_1,x_2){=}0$, $f(x_1,x_2){=}x_1^{-1/2}$\linebreak ($x_1\ne0$). 


 


Тогда $\rho(f,H)=0$. Нетрудно видеть, что здесь нет элемента наилучшего 


несимметричного приближения. 


\end{exam} 


 


Для доказательства теорем понадобятся некоторые леммы. Следуя [3] 


и [4], введем дополнительные понятия и обозначения. 


 


Пусть $H$ --- конечномерное подпространство в $L_0(G)$, 


$Q=\{q_i\}^m_{i=1}$ --- его базис. Для $g\in H$ положим 


$\|g\|_Q=\sum\limits^m_{i=1}|\lambda_i|$,


где $g(x)=\sum\limits^m_{i=1}\lambda_i 


q_i(x)$. Очевидно, что введенная величина является нормой на $H$. Назовем 


ее нормой, порожденной базисом $Q$. Множество $T\subset H$ будем называть 


$\lambda$-ограниченным, если $\sup\limits_{g\in T}\|g\|_Q<\infty$. Ясно, 


что понятие $\lambda$-ограниченности инвариантно относительно выбора 


базиса. Всякое $\lambda$-ограниченное множество $T\subset H$, очевидно, 


компактно в смысле сходимости п.\,в. (т.\,е. из любой последовательности 


$\{g_i\}^\infty_{i=1}$ можно выбрать подпоследовательность, сходящуюся 


п.\,в. к некоторой функции из $H$). Если $\|g\|_Q=1$, то элемент $g$ 


назовем нормированным (относительно выбранного базиса). 


 


Последовательность элементов $\{g_k\}^\infty_{k=1}\subset H$ назовем 


минимизирующей для $f\in L_{\varphi,G}$, если $\rho(f,g_k)<\infty$ при 


всех $k=1,2,\dotsc$ и $\lim\limits_{k\to\infty}\rho(f,g_k)=\rho(f,H)$. 


 


 


\begin{lem} 


Пусть $Q=\{g_i\}^m_{i=1}$ --- базис в $H\subset L_{\varphi,G}$, а функция 


$p\in L_0(G)$ такова, что $\mu\{x\in G:p(x)=0\}=0$. 


 


Тогда найдется такое натуральное число $M$, что на множестве 


$G_M=\{x\in G:|p(x)|\ge\frac{1}{M}\}$ \ $\{q_i\}^m_{i=1}$ будет базисом 


$\Big($т.\,е. $\sum\limits^m_{i=1}\lambda_i q_i(x)=0$ 


для п.\,в. $x\in G_M$ лишь при $\lambda_1=\dotsb=\lambda_m=0\Big)$. 


\end{lem} 


 


 


\begin{pf} 


Допустим, что лемма неверна. Тогда на каждом из множеств 


последовательности $G_1\subset G_2\subset \dotsb \subset G_k\subset 


\dotsb \subset G$ система $Q$ не будет базисом, причем $\mu(G\setminus 


G_0)=0$ для $G_0=\bigcup\limits^\infty_{k=1}G_k$. Поскольку 


$\{q_i\}^m_{i=1}$ на $G_k$ --- не базис, то найдутся числа 


$\lambda^{(k)}_1,\dotsc,\lambda^{(k)}_m$ такие, что 


$\Lambda_k=|\lambda_1^{(k)}|+ \dotsb+ |\lambda^{(k)}_m|\ne0$ и 


$\lambda^{(k)}_1q_1(x)+ \dotsb+ \lambda^{(k)}_m q_m(x)=0$ п.\,в. на 


$G_k$. Так как $\Lambda_k\ne0$, то $g_k(x)=(\lambda_1^{(k)}q_1(x)+ 


\dotsb+ \lambda_m^{(k)}q_m(x))/\Lambda_k$ также равна п.\,в. на $G_k$ 


нулю. В силу $\lambda$-огра\-ни\-ченности множества $\{g_k\}^\infty_{k=1}$ 


(поскольку $\|g_k\|_Q=1$) можно выделить подпоследовательность 


$\{g_{k_s}\}$, сходящуюся п.\,в. к нормированному (а значит, 


нетривиальному) элементу $g\in H$. Тогда п.\,в. 


$\lim\limits_{s\to\infty}g_{k_s}(x)\chi_{G_{k_s}}(x)= g(x)\chi_{G_0}(x)$, 


где $\chi_{G_k}(x)$ --- характеристическая функция множества $G_k$. С 


другой стороны, $g(x)$ п.\,в. равна нулю на $G$, что невозможно. 


\end{pf} 


 


 


\begin{lem} 


Пусть $Q=\{q_i\}^m_{i=1}$ --- базис в $H\subset L_{\varphi,G}$. Тогда для 


любого измеримого множества $T\subset G$, на котором система $Q$ является 


базисом, найдутся такие числа $\varepsilon_0>0$ и $\delta_0>0$, что для 


любого $g\in H$ 


$$ 


\mu\{x\in T:|g(x)|\ge\|g\|_Q\delta_0\}\ge\varepsilon_0. 


$$ 


\end{lem} 


 


 


\begin{pf} 


Допустим, что лемма неверна. Тогда существует последовательность таких 


нормированных элементов $\{g_k\}$, что 


\begin{equation} 


\mu\{x\in G:|g_k(x)|\ge1/k\}\le 1/k \ \ (k=1,2,\dotsc). 


\end{equation} 


 


В силу компактности множества нормированных элементов из 


последовательности $\{g_k\}$ можно выделить подпоследовательность 


$\{g_{k_s}\}$, сходящуюся п.\,в. к нормированному (а значит, 


нетривиальному) элементу $g_0$. Но в силу (5) $\{g_{k_s}\}$ сходится по 


мере к тождественному нулю. Отсюда вытекает $g_0(x)=0$ п.\,в., что 


невозможно. 


\end{pf} 


 


 


\begin{lem}[{\series{m}\shape{n}\selectfont [5], с.\,32}] 


Если  последовательность измеримых неотрицательных функций $\{f_k\}$


на измеримом множестве $G\subset R^n$ сходится к функции $f(x)$, то 


$$ 


\int_G f(x)dx\le\sup\bigg\{\int_G f_k(x)dx\bigg\}. 


$$ 


\end{lem} 


 


 


\begin{lem}[{\series{m}\shape{n}\selectfont [6], с.\,13}] 


Если последовательность измеримых функций $\{f_k\}$ сходится к $f$ на 


множестве $G$ и $|f_k(x)|\le\varphi(x)$ для некоторой интегрируемой на 


$G$ функции $\varphi(x)$ при любом $k$, то предельная функция 


интегрируема на $G$ и 


$$ 


\int_G f_k(x)dx\to\int_G f(x)dx\quad (k\to\infty). 


$$ 


\end{lem} 


 


 


\begin{lem}[{\series{m}\shape{n}\selectfont [7], с.\,131}] 


Непрерывная на оси $Ox$ выпуклая вниз функция $M(x)$ имеет в каждой точке 


правую $M'_\sqcap(x)$ и левую $M'_\wedge(x)$ производные, причем 


$M'_\wedge(x)\le M'_\sqcap(x)$. Если $a<x_1<x_2<b$, то 


$M'_\sqcap(a)\le(M(x_2)-M(x_1))/ (x_2-x_1)\le M'_\wedge(b)$. 


\end{lem} 


 


 


\begin{lem} 


Для любых $f,g,p=(p_1,p_2)$ из $L_0(G)$ справедливо 


неравенство $($при п.\,в. $x\in G)$ 


$$ 


|((f(x)+g(x));p(x))|\le|(f(x);p(x))|+|(g(x);p(x))|. 


$$ 


\end{lem} 


 


{\bf Доказательство} проводится непосредственной проверкой. 


 


 


\begin{lem} 


Пусть $\varphi\in\phi$, $H$ --- конечномерное подпространство в $L_{\varphi,G}$, 


$p=(p_1,p_2)$ --- такой вес, что $gp_1\in L_{\varphi,G}$, $gp_2\in 


L_{\varphi,G}$ для любого элемента $g\in H$. 


 


Для того чтобы функция $g_0\in H$ была элементом наилучшего 


несимметричного приближения для функции $f\in L_{\varphi,G}$, для которой 


$(f;p)\in L_{\varphi,G}$, необходимо, чтобы для любой функции $g\in H$ 


выполнялось условие $(3)$. 


\end{lem} 


 


 


\begin{pf} 


Считаем, что $\varphi(u)\not\equiv\const$, т.\,к. иначе лемма 7 


тривиальна. 


 


Пусть $Q=\{q_i\}^m_{i=1}$ --- базис конечномерного подпространства $H$, 


функция $f\in L_{\varphi,G}$. Ниже будем рассматривать только те точки 


$x$, в которых конечны функции $f(x)$, $q_i(x)$, $i=1,\dotsc,m$, т.\,к. 


мера точек, не удовлетворяющих этому условию, равна нулю. 


 


Пусть $g_0(x)$ является элементом наилучшего несимметричного приближения, 


$g(x)$ --- произвольный элемент из $H$. Тогда в базисе $Q$ их можно 


представить в виде $g_0(x)=\sum\limits^m_{i=1}a_i^{(0)}q_i(x)$, 


$g(x)=\sum\limits^m_{i=1}a_i q_i(x)$. Рассмотрим функцию 


$$ 


I(a_1,\dotsc,a_m)=\int_G \varphi\bigg(\bigg|\bigg( 


\bigg(f(x)-\sum^m_{k=1} 


a_k q_k(x)\bigg);p(x)\bigg)\bigg|\bigg)dx. 


$$ 


 


Возьмем какую-либо последовательность $\{h_s\}^\infty_{s=1}$, для которой 


$0<h_s\le1$, $h_s\to0$ ($s\to\infty$), и найдем предел 


$$ 


\lim_{s\to+\infty}(I(a_1^{(0)}+h_s(a_1-a_1^{(0)}),\dotsc, 


a_m^{(0)}+h_s(a_m-a_m^{(0)}))-I(a_1^{(0)},\dotsc,a_m^{(0)})) 


/h_s=\lim_{s\to+\infty}\int_G I_s(x)dx, 


$$ 


где 


$$ 


I_s(x)=\{\varphi(|((F(x)-h_s G(x));p(x))|)- 


\varphi(|(F(x);p(x))|)\}/h_s,\ \ 


F(x)=f(x)-g_0(x),\ \ G(x)=g(x)-g_0(x). 


$$ 


 


Выясним, к чему сходятся подинтегральные функции. Для этого совершим 


предельный переход в трех случаях: 


 


1) $F(x)>0$, \ 2) $F(x)<0$, \ 3) $F(x)=0$. 


 


В случае 1), если $s$ достаточно велико, 


$$ 


I_s=\{\varphi((F(x)-h_s G(x))|p_1(x)|)- 


\varphi(F(x)|p_1(x)|)\}/h_s. 


$$ 


Если $G(x)>0$, то при $s\to\infty$ 


\begin{multline*} 


I_s\to\varphi'_\wedge(F(x)|p_1(x)|)(-|p_1(x)|G(x))= 


-\{\varphi'_\wedge(|(F(x);p(x))|) \sign\{F(x)G(x)\}^+ +\\ 


% 


+\varphi'_\sqcap(|(F(x);p(x))|)\sign\{F(x)G(x)\}^-\} 


G(x)\{|p_1(x)|\sign F^+(x)-|p_2(x)|\sign F^-(x)\}=B(x). 


\end{multline*} 


 


Если $G(x)<0$, то при $s\to\infty$ \ $I_s\to\varphi'_\sqcap 


(F(x)|p_1(x)|)(-|p_1(x)|G(x))=B(x)$. 


Если $G(x)=0$, то $I_s=0=B(x)$. 


 


Случай 2) аналогичен разобранному. 


 


Рассмотрим теперь случай 3). Если $G(x)>0$, то $I_s=(\varphi(|p_2(x)|h_s 


G(x))-\varphi(0))/h_s$ при достаточно больших $s$. Тогда при $s\to\infty$ 


будем иметь $I_s\to \varphi'_\sqcap(0) |p_2(x)|G(x)= 


\varphi'_\sqcap(0)|(G;p)(x)|$. 


 


Если $G(x)<0$, то $I_s=(\varphi(-|p_1(x)|h_s(x)G(x))- \varphi(0))/h_s$ 


при достаточно больших $s$ и при $s\to\infty$ будем иметь 


$I_s\to-\varphi'_\sqcap (0)|p_1(x)|G(x)= \varphi'_\sqcap(0) |(G;p)(x)|$. 


 


Сделаем общий вывод. Подинтегральные функции при $s\to\infty$ стремятся к 


функции $F_1(x)$, определенной следующим образом: при $x\in G\setminus E$ 


функция $F_1(x)$ равна $B(x)$, а при $x\in E$ --- 


$\varphi'_\sqcap(0)|(G;p)(x)|$. 


 


Покажем теперь, что для любого $s$ подинтегральные функции по модулю 


не превосходят $cB_1(x)$, где 


$$ 


B_1(x)=\varphi(|(f(x);p(x))|)+\varphi(|g(x)p_1(x)|) 


+\varphi(|g_0(x)p_1(x)|)+\varphi(|g(x)p_2(x)|)+ 


\varphi(|g_0(x)p_2(x)|),


$$ 


$c$ --- некоторая фиксированная постоянная, зависящая от $\varphi$. 


 


Пусть $G(x)>0$. Если $F(x)>0$, то разность $R(x)=F(x)-h_s G(x)$ 


может быть положительной, отрицательной или равной нулю. В точках, 


где она положительна, умножив и разделив $|I_s|$ на величину 


$|p_1(x)|G(x)$ и применяя леммы 5, 6, имеем 


$$ 


|I_s|\le \varphi'_\wedge(|p_1(x)|F(x))| 


p_1(x)G(x)|\le c_1B_1(x). 


$$ 


В точках, где $R(x)=0$, вновь применяя леммы 5 и 6, получим 


$$ 


|I_s|=\frac{1}{h_s}|\varphi(0)-\varphi(|p_1(x)|F(x))|\le 


\varphi'_\wedge(|p_1(x)|F(x))|p_1(x)|G(x)\le c_2B_1(x). 


$$ 


В точках, где $R(x)<0$, будем иметь 


\begin{multline*} 


|I_s|=\frac{1}{h_s}|\varphi(-|p_2(x)|(F(x)-h_s G(x)))- 


\varphi(|p_1(x)|F(x))|\le \\ 


% 


\le\frac{1}{h_s}|\varphi(-|p_2(x)|(F(x)-h_s G(x)))- 


\varphi(0)|+\frac{1}{h_s}|\varphi(|p_1(x)|F(x))- 


\varphi(0)|=I^{(1)}_s+I^{(2)}_s. 


\end{multline*} 


Оценим $I^{(1)}_s$ аналогично тому, как в случае $R(x)>0$: 


$$ 


I^{(1)}_2\le \varphi'_\wedge(-|p_2(x)|(F(x)-h_s G(x))) 


(F(x)-h_s G(x))(-|p_2(x)|)/h_s\le c_3B_1(x) 


$$ 


(т.\,к. $-(F(x)-h_s G(x))\le h_s G(x)$). Аналогично 


$$ 


I^{(2)}_2\le c_4B_1(x) 


$$ 


(т.\,к. $F(x)=F(x)-h_s G(x)+h_s G(x)\le h_s G(x)$). Объединяя оценки для 


$I^{(1)}_s$ и $I^{(2)}_s$, будем иметь $|I_s|\le c_5B_1(x)$. Аналогично 


получим, что $|I_s|\le c_6B_1(x)$ и в случае, когда $F(x)\le0$. Итак, 


случай, когда $G(x)>0$, разобран. 


 


Аналогичными рассуждениями получаем верхнюю оценку для $|I_s|$ через 


$B_1(x)$ в случае, когда $G(x)<0$. Если же $G(x)=0$, то оценка сверху 


тривиальна, т.\,к. $|I_s|=0\le B_1(x)$. 


 


Таким образом, показано, что $|I_s|\le cB_1(x)$, где $B_1(x)$ является 


суммируемой функцией, для любого $s$. Значит, последовательность 


$\{I_s\}$ удовлетворяет всем условиям леммы 4, применение которой 


позволяет совершить предельный переход под знаком интеграла. А так как 


$g_0(x)$ является элементом наилучшего несимметричного приближения, то 


$$ 


I(a_1^{(0)}+h_s(a_1-a_1^{(0)}),\dotsc,a_m^{(0)}+ 


h_s(a_m-a_m^{(0)}))-I(a_1^{(0)},\dotsc,a_m^{(0)})\ge0,


$$ 


и поэтому 


\begin{equation} 


\int_G F_1(x)dx\ge0. 


\end{equation} 


После замены функции $G(x)$, входящей в выражение для $F_1(x)$, на 


$-g(x)$ неравенство (6) перейдет в (3). 


\end{pf} 


 


 


\begin{lem} 


Пусть $\varphi\in\phi$, $p=(p_1,p_2)$ --- вес, $H$ --- конечномерное 


подпространство в $L_{\varphi,G}$, $f\in L_{\varphi,G}$, $g_0\in H$, 


$((f-g_0);p)\in L_{\varphi,G}$. 


 


Если для любой функции $g\in H$ выполняется условие $(3)$, то 


функция $g_0$ является элементом наилучшего несимметричного приближения. 


\end{lem} 


 


 


{\bf Доказательство.} Пусть для любого $g\in H$ выполняется 


(3), а значит, и (6). Справедливы следующие неравенства: 


\begin{multline*} 


M=\int_{G\setminus E}\{\varphi'_\wedge(|(F(x);p(x))|) 


\sign\{F(x)G(x)\}^+ +\\ 


% 


+\varphi'_\sqcap(|(F(x);p(x))|)\sign\{F(x)G(x)\}^-\} 


G(x)(\sign F(x)|p_1(x)|-\sign F^-(x)|p_2(x)|)dx= \\ 


% 


=\int_G(-F_1(x))dx+\int_E F_1(x)dx+\int_{G\setminus E} 


\{\varphi'_\wedge(|(F(x);p(x))|) 


\sign\{F(x)G(x)\}^+ +\\ 


% 


+\varphi'_\sqcap(|(F(x);p(x))|)\sign\{F(x)G(x)\}^-\} 


G(x)(\sign F(x)|p_1(x)|-\sign F^-(x)|p_2(x)|)dx, 


\end{multline*} 


где $F(x)=f(x)-g_0(x)$, $G(x)=g(x)-g_0(x)$. Учитывая (6) и проводя 


преобразования в третьем интеграле, получим 


\begin{multline*} 


M\le\int_E F_1(x)dx+\int_{G\setminus E} 


\{\varphi'_\wedge(|(F(x);p(x))|)\sign 


\{F(x)G(x)\}^+ + \\ 


% 


+\varphi'_\sqcap(|(F(x);p(x))|)\sign\{F(x)G(x)\}^-\} 


(|((f-g);p)(x)|-|(F;p)(x)|)dx+M. 


\end{multline*} 


Отсюда имеем 


\begin{multline*} 


B_2=\int_E F_1(x)dx+\int_{G\setminus E} 


\{\varphi'_\wedge(|(F(x);p(x))|)\sign\{F(x)G(x)\}^+ +\\ 


% 


+\varphi'_\sqcap(|(F(x);p(x))|)\sign\{F(x)G(x)\}^-\} 


(|((f-g);p)(x)|-|(F;p)(x)|)dx\ge0. 


\end{multline*} 


Используя лемму 5, для $\varphi\in\phi$ получим 


\begin{equation} 


\varphi(|a+b|)\ge\varphi(|a|)+b\sign a 


\{\varphi'_\wedge(|a|)\sign B^+ +\varphi'_\sqcap(|a|) 


\sign B^-\} 


\end{equation} 


при любом $B\ne0$. Если $b=0$, то неравенство (7) будет справедливо и 


для $B=0$. Применяя (7), для любого $g\in H$ будем иметь 


\begin{multline*} 


\int_G \varphi(|((f-g);p)|)dx=\int_G \varphi(|(F;p)|+ 


(|((f-g);p)|-|(F;p)|))dx\ge \\ 


% 


\ge B_2+\int_{G\setminus E} \varphi(|(F;p)|)dx+ 


\int_E \varphi(0)dx\ge\int_G \varphi(|(F;p)|)dx.\quad\square 


\end{multline*} 


 


 


\begin{pf*}{Доказательство теоремы 1} 


Если $\rho(f,H)=\infty$ или функция $\varphi$ --- константа, то 


утверждение теоремы 1 очевидно. Пусть $\rho(f,H)<\infty$, 


$\varphi\in\phi$ и отлична от константы. Тогда 


$\lim\limits_{u\to+\infty}\varphi(u)=+\infty$. 


 


Пусть $\{g_k\}$ --- минимизирующая для $f$ последовательность элементов 


из $H$. Начиная с некоторого $k_0$, для всех $k$ справедливо неравенство 


$\rho(f,g_k)\le\rho(f,g_1)+1$. Тогда найдется такое положительное число 


$A$, что $\|g_k\|_Q\le A<\infty$ для всех $k=k_0+1,\dotsc$ Докажем это 


методом от противного. Если это не так, то для любого положительного 


числа $B$ найдется такое число $k\ge k_0+1$, что $\|g_k\|_Q>B$. В силу 


леммы 1 найдутся натуральное число $m_0$ и множество 


$G^{(1)}_{m_0}=\{x\in G:|p_1(x)|\ge 1/m_0\}$, на котором система 


$Q$ является базисом. Также в силу леммы 1 найдутся натуральное число 


$s_0$ и множество $G^{(2)}_{s_0}=\{x\in G^{(1)}_{m_0}:|p_2(x)|\ge 


1/s_0\}$, причем система $Q$ будет на нем базисом. Тогда на множестве 


$G_0=G^{(2)}_{s_0}$ будем иметь $|p_1(x)|\ge 1/m_0$, $|p_2(x)|\ge 1/s_0$. 


Из леммы 2 следует существование таких $\varepsilon_0>0$ и $\delta_0>0$, 


что для любого $g\in H$ мера множества $R_g=\{x\in G_0:|g(x)|\ge 


\|g\|_Q\delta_0\}$ больше $\varepsilon_0$. Пусть $R^N_g=R_g\cap \{x\in 


G:|f(x)|\le N\}$. Функция~$f$ конечна п.\,в., поэтому найдется такое 


$N_0$, что при $N>N_0$ и всех $g\in H$ (в частности, и для минимизирующей 


последовательности $\{g_k\}$) будет справедливо неравенство 


$\mu(R^N_g)>\varepsilon_0/2$. Тогда $\rho(f,g_k)\ge\int\limits_{R^N_{g_k}} 


\varphi(|(F_k(x);p(x))|)dx$, где $F_k(x)=f(x)-g_k(x)$. 


 


Рассмотрим множества 


$$ 


G^+_k=\{t\in R^N_{g_k}:F_k(x)>0\},\quad 


G^-_k=\{t\in R^N_{g_k}:F_k(x)\le 0\}. 


$$ 


Заметим, что $G^+_k\cup G^-_k=R^N_{g_k}$, $G^+_k\cap G^-_k=\emptyset$. 


Тогда


$$


F_{k_s}(x)p(x)\rho(f,g_k)\ge c_1\bigg(\int_{G^+_k} \varphi(F_k(x))dx+ 


\int_{G^-_k} \varphi(-F_k(x))dx\bigg),


$$


где $c_1$ --- положительная 


постоянная, зависящая лишь от функции $\varphi$ и чисел $m_0$, $s_0$. 


 


Возьмем $N$ настолько большим, что $\varphi(N)\ge2(\rho(f,g_1)+1)/ 


(\varepsilon_0c_1)$. Тогда для тех элементов $g_k\in H$, для которых 


$\|g_k\|_Q>B_0=2N/\delta_0$, выполняется неравенство 


$\rho(f,g_k)>\rho(f,g_1)+1$, что противоречит ограничению $\rho(f,g_k)\le 


\rho(f,g_1)+1$. Полученное противоречие доказывает, что $\{g_k\}$ 


является $\lambda$-ограниченной последовательностью. 


 


Так как последовательность $\{g_k\}$ \ $\lambda$-ограничена, то 


существуют такие последовательность $\{g_{k_s}\}$ и элемент $g\in H$, что 


$\lim\limits_{s\to\infty}g_{k_s}(x)=g(x)$ при п.\,в. $x\in G$. Из леммы 3 


следует, что $\rho(f,g)\le \lim\limits_{s\to\infty} \int\limits_G 


\varphi(|(F_{k_s}(x);p(x))|)dx= \rho(f,H)$, а значит, 


$\rho(f,g)=\rho(f,H)$, т.\,е. $g$ является элементом наилучшего 


несимметричного приближения. 


 


Покажем теперь, что если функция $\varphi$ дополнительно удовлетворяет 


условию строгой выпуклости, то элемент наилучшего несимметричного 


приближения является единственным. Доказательство проведем от противного. 


Пусть это не так, т.\,е. найдется хотя бы еще одна функция $g_1(x)\in H$, 


отличная от $g(x)$ и являющаяся элементом наилучшего несимметричного 


приближения функции $f$. Тогда для функции $(g(x)+g_1(x))/2$, используя 


свойства функции $\varphi$, будем иметь 


\begin{multline*} 


\rho(f,(g+g_1)/2)=\int_G \varphi(|((f(x)-(g(x)+g_1(x))/2); 


p(x))|)dx < \\ 


% 


<\bigg(\int_G \varphi(|((f(x)-g(x));p(x))|)dx+ 


\int_G \varphi(|((f(x)-g_1(x));p(x))|)dx\bigg)/2=\rho(f,H), 


\end{multline*} 


что невозможно. Таким образом, элемент $g$ является единственным 


элементом наилучшего несимметричного приближения функции $f$. 


\end{pf*} 


 


 


\begin{pf*}{Доказательство теоремы 2} 


Из условий теоремы 2 следует, что условия лемм 7 и 8 выполнены. 


Применение этих лемм позволяет получить справедливость теоремы 2. 


\end{pf*} 


 


 


\begin{thebibliography}{9} 


 


\bibitem{1} 


Бабенко В.Ф. {\em Несимметричные приближения в пространствах суммируемых 


функций} // Укр. матем. журн. -- 1982. -- Т.\,34. 


-- \No\,4. -- С.\,409--419. 


\bibitem{2} 


Корнейчук Н.П. {\em Экстремальные задачи теории приближения}. -- М.: 


Наука, 1978. -- 320 с. 


\bibitem{3} 


Гаркави А.Л. {\em Теорема существования элемента наилучшего приближения 


в пространствах типа $(F)$ с интегральной метрикой} // Матем. заметки. -- 


1970. -- Т.\,8. -- \No\,5. -- С.\,583--594. 


\bibitem{4} 


Рахметов Н.К. {\em О некоторых вопросах теории приближения функций}\/: 


Дисс.~$\dotsc$ канд. физ.-матем. наук. -- М.: МГУ, 1990. -- 125~с. 


\bibitem{5} 


Никольский С.М. {\em Приближение функций многих переменных и теоремы 


вложения}. -- М.: Наука, 1975. -- 455 с. 


\bibitem{6} 


Бесов О.В., Ильин В.П., Никольский С.М. {\em Интегральные представления 


функций и теоремы вложения}. -- М.: Наука, 1977. -- 480 с. 


\bibitem{7} 


Красносельский М.А., Рутицкий Я.Б. {\em Выпуклые функции и пространства 


Орлича}. -- М.: Физматгиз, 1958. -- 271 с. 


 


\end{thebibliography} 


 


\vskip 2cm 


 


\post{\it Волгоградский государственный}{\it Поступила} 


\post{\it технический университет}{28.04.2003} 


 


\label{end} 


\end{document} 


 


